\section{The inherited physical and coherent interface}
\label{sec:tpc37-inherited-interface}

We restate only the inputs used below.  They come from the actual
source-compatible packet and the reductions recorded in TPC-25,
TPC-35, and TPC-36 \citep{WangTPC25,WangTPC35,WangTPC36}.  Finite-field
identities proved later in this paper are unconditional algebraic
statements, but their application to the prime--M\"obius packet is always
made through the interface in this section.

\subsection{Rows, orbit, and ultra increment}

Fix $h\in\mathbb Z\setminus\{0\}$, put $H=\operatorname{rad}|h|$,
and work on one dyadic packet with
\begin{equation}
 Q=LD,\qquad QJ\asymp X,\qquad R\le T<U_0\asymp X.
 \label{eq:tpc37-interface-scales}
\end{equation}
A physical row is
\begin{equation}
 \alpha=(\ell_\alpha,d_\alpha),\qquad
 m_\alpha=\ell_\alpha d_\alpha\asymp Q,
 \label{eq:tpc37-interface-row}
\end{equation}
where $\ell_\alpha\asymp L$ is prime,
$d_\alpha\asymp D$ is squarefree, and
$(d_\alpha,H)=1$.  The source separation makes the integer $m_\alpha$
determine its row uniquely.  There are $O(Q)$ rows.  For the two
physical coefficient systems,
\begin{align}
 \sum_\alpha|\gamma_\alpha^{(i)}|
 &\ll X^{o(1)}Q,
 \label{eq:tpc37-interface-row-l1}\\
 \sum_\alpha|\gamma_\alpha^{(i)}|^2
 &\ll X^{o(1)}Q,
 \qquad i=1,2.
 \label{eq:tpc37-interface-row-l2}
\end{align}
These estimates retain the M\"obius/source provenance used in the
preceding papers; they are not hypotheses on an arbitrary $Q$-entry
vector.

The orbit variable $j$ lies in a positive interval
$\mathcal J$ containing $O(J)$ integers, and
\begin{equation}
 N_\alpha(j):=m_\alpha j+h\asymp X.
 \label{eq:tpc37-interface-target}
\end{equation}
The terminal range is chosen so that
\begin{equation}
 T<\min_{\alpha,j}N_\alpha(j)
 \le\max_{\alpha,j}N_\alpha(j)<U_0.
 \label{eq:tpc37-interface-terminal-range}
\end{equation}
Thus the terminal divisor $u=N_\alpha(j)$ is present in the ultra
increment
\begin{equation}
 C_m(j)=
 \sum_{\substack{T<u\le U_0\\u\mid mj+h}}
 -\mu(u)\log u,
 \qquad |C_m(j)|\ll X^{o(1)}.
 \label{eq:tpc37-interface-ultra-increment}
\end{equation}
On the physical support every divisor $u\mid N_\alpha(j)$ obeys
\begin{equation}
 (u,m_\alpha jH)=1.
 \label{eq:tpc37-interface-primitivity}
\end{equation}

The actual multiplier has the form
\begin{equation}
 \mathfrak A_{\alpha,\gamma}^{\rm act}(j)
 =\mathfrak m_{\rm act}(\alpha,\gamma)
  \Xi_{\alpha,\gamma}(j)
  \mathcal W_{\alpha,\gamma}(j/J),
 \label{eq:tpc37-interface-multiplier}
\end{equation}
where all displayed periodic and smooth factors are bounded and
\begin{equation}
 \mathfrak m_{\rm act}(\alpha,\gamma)
 =\mathbf1_{\ell_\alpha\ne\ell_\gamma}
  \mathbf1_{|m_\alpha-m_\gamma|>\Delta}
  \mathbf1_{(d_\alpha,d_\gamma)\le G}
 \label{eq:tpc37-interface-static-mask}
\end{equation}
for the inherited pruning scales
$\Delta=QX^{-\kappa_0}$ and $G=X^{\kappa_0}$.
TPC-36 proves that this static mask, together with the fixed orbit cells
and smooth factors, has $X^{o(1)}$ signed product-ready reassembly cost.
That theorem does not complete the integer ranges modulo $q$, nor does
it estimate the resulting coherent divisor sum.

\subsection{The two orbit energies and their gate}

Define the left slice and its energy by
\begin{align}
 Y_{\gamma,j}^L
 &:=\sum_\alpha \gamma_\alpha^{(1)}
 \mathfrak A_{\alpha,\gamma}^{\rm act}(j)C_{m_\alpha}(j),
 \label{eq:tpc37-interface-left-slice}\\
 V_L&:=\sum_{\gamma,j}|Y_{\gamma,j}^L|^2.
 \label{eq:tpc37-interface-left-energy}
\end{align}
The right slice exchanges the two row systems:
\begin{align}
 Y_{\alpha,j}^R
 &:=\sum_\gamma \gamma_\gamma^{(2)}
 \mathfrak A_{\alpha,\gamma}^{\rm act}(j)C_{m_\gamma}(j),
 \label{eq:tpc37-interface-right-slice}\\
 V_R&:=\sum_{\alpha,j}|Y_{\alpha,j}^R|^2.
 \label{eq:tpc37-interface-right-energy}
\end{align}
The inherited scalar residual is controlled once
\begin{equation}
 V_L+V_R\ll X^{o(1)}\frac{Q^3}{J}
 \label{eq:tpc37-interface-orbit-gate}
\end{equation}
is proved on every retained packet.  This is a sufficient interface, not
a conclusion imported into this paper.  The complete input-copy diagonal
has already been bounded by
\begin{equation}
 V_{L,\Delta}+V_{R,\Delta}
 \ll X^{o(1)}Q^2J.
 \label{eq:tpc37-interface-diagonal}
\end{equation}
Consequently the outstanding task is coefficient-specific and
off-diagonal; replacing it by a maximum over all abstract row vectors is
unnecessarily strong.

For later sign-blind comparisons, put
\begin{equation}
 C_m^{\rm abs}(j)
 :=\sum_{\substack{T<u\le U_0\\u\mid mj+h}}\log u.
 \label{eq:tpc37-interface-absolute-increment}
\end{equation}
The divisor bound gives $C_m^{\rm abs}(j)\ll X^{o(1)}$.  Define
$Y^{L,\rm abs}$ from \eqref{eq:tpc37-interface-left-slice} by replacing
$\gamma_\alpha^{(1)}$,
$\mathfrak A_{\alpha,\gamma}^{\rm act}(j)$, and the ultra summands by
their absolute values, and define $Y^{R,\rm abs}$ symmetrically.  Then
the row $\ell^1$-bound, the $O(Q)$ opposite rows, and the $O(J)$
orbit times give
\begin{equation}
 \boxed{
 V_{\rm eq}^{\rm abs}
 :=\sum_{\gamma,j}|Y_{\gamma,j}^{L,\rm abs}|^2
  +\sum_{\alpha,j}|Y_{\alpha,j}^{R,\rm abs}|^2
 \ll X^{o(1)}Q^3J.}
 \label{eq:tpc37-interface-absolute-baseline}
\end{equation}
This estimate is deliberately crude: it records the exact-equality
envelope after all M\"obius signs have been erased.  It is used only when
an algebraic face has a uniform scalar amplitude on $su=N_\alpha(j)$.
It is not evidence of cancellation and cannot by itself imply
\eqref{eq:tpc37-interface-orbit-gate}.

\subsection{Reciprocal dyadic blocks and product slopes}

Open one divisor in \eqref{eq:tpc37-interface-ultra-increment} and write
\begin{equation}
 su=N_\alpha(j)=\ell_\alpha d_\alpha j+h.
 \label{eq:tpc37-interface-reflection}
\end{equation}
Partition $u$ into $O(\log X)$ reciprocal dyadic blocks
$u\asymp U$, $U\ge T$, and attach the corresponding interval
\begin{equation}
 s\asymp X/U.
 \label{eq:tpc37-interface-complementary-range}
\end{equation}
On the full Cartesian support of each block, $su=O(X)$.  After fixing
the opposite row, one orbit residue cell, and one bounded projective
layer from TPC-36, the integer core has the schematic product-ready form
\begin{equation}
 \sum_{\ell,d,u,s}
 b(\ell)f(d)g(u)w(s)
 \mathbf1_{su=\ell jd+h}.
 \label{eq:tpc37-interface-integer-core}
\end{equation}
Here $b$ is independent of $j$ and $s$, and the physical ultra
coefficient on a clean dyadic block is
\begin{equation}
 g_{\rm phys}(u)=-\mu(u)\log u\,\Phi_U(u)
 \label{eq:tpc37-interface-physical-ultra}
\end{equation}
with a controlled dyadic cutoff $\Phi_U$.  The weight $w$ records the
incomplete complementary interval and any alias multiplicity introduced
later.  Formula \eqref{eq:tpc37-interface-integer-core} is a topology and
provenance statement; it is not a complete finite-field identity.

For a prime $q\nmid h$, the corresponding complete product--slope model
is defined on $G_q=\mathbb F_q^\times$ by
\begin{align}
 R_s(a)
 &:=\sum_{D\in\mathbb F_q}
 f(D)g\!\left(s^{-1}(aD+h)\right),
 \label{eq:tpc37-interface-Rsa}\\
 (Tb)(j,s)
 &:=\sum_{\ell\in G_q}b(\ell)R_s(\ell j),
 \label{eq:tpc37-interface-Tb}\\
 S_w(j)&:=\sum_{s\in G_q}w(s)(Tb)(j,s).
 \label{eq:tpc37-interface-Sw}
\end{align}
With the unnormalized multiplicative transforms
\begin{equation}
 \widehat b(\chi)=\sum_{x\in G_q}b(x)\overline{\chi(x)},
 \qquad
 \widehat R_s(\chi)=\sum_{a\in G_q}R_s(a)\overline{\chi(a)},
 \label{eq:tpc37-interface-transforms}
\end{equation}
TPC-36 gives the exact coherent identity
\begin{equation}
 \boxed{
 \mathcal G_q(b,f,g;w)
 :=\sum_{j\in G_q}|S_w(j)|^2
 =\frac1{q-1}\sum_\chi
 |\widehat b(\bar\chi)|^2
 \left|\sum_{s\in G_q}w(s)\widehat R_s(\chi)\right|^2.}
 \label{eq:tpc37-interface-coherent-gate}
\end{equation}
The right side is the actual finite-model coherent functional, not an
upper bound.  The square-function spectrum
$\sum_{j,s}|(Tb)(j,s)|^2$ controls an average over independent $s$'s;
it does not control \eqref{eq:tpc37-interface-coherent-gate} without a
possible factor $q-1$.  TPC-36 proves that this factor is sharp for
arbitrary centered data.

\subsection{Endpoint and alias ledger}

At the endpoint used throughout this paper,
\begin{equation}
 Q=X^{267/400+o(1)},\qquad
 J=X^{133/400+o(1)},\qquad
 T=X^{193/500+o(1)},
 \label{eq:tpc37-interface-endpoint}
\end{equation}
and the source scales satisfy $D=o(J)$, $J=o(L)$, and
$QJ=X^{1+o(1)}$.  The identities relevant to the present ledger are
\begin{align}
 \frac Q{J^2}
 &=X^{1/400+o(1)},
 &\frac X{J^3}
 &=X^{1/400+o(1)},
 \label{eq:tpc37-interface-critical-identities}\\
 \frac T J
 &=X^{107/2000+o(1)},
 &\frac QT
 &=X^{563/2000+o(1)}.
 \label{eq:tpc37-interface-secondary-identities}
\end{align}
One modulus $q\asymp J$ leaves $O(1+X/q)=O(Q)$ possible integer
determinants in a block.  Three pairwise coprime orbit-scale moduli leave
\begin{equation}
 O\!\left(1+\frac X{J^3}\right)
 =X^{1/400+o(1)}
 \label{eq:tpc37-interface-three-modulus-aliases}
\end{equation}
joint aliases.  This count matches the allowance above the diagonal, but
it does not show that principal, proper-coordinate, or full-coordinate
faces reassemble with the required sign.

The logical division used below is therefore strict:
\begin{enumerate}[label=(\roman*)]
 \item the physical equality is used first to remove residue-zero
 coordinates;
 \item punctured centering and two-character analysis concern the
 regular complete finite model;
 \item the three-modulus face calculus controls exact periodic masses and
 determinant multiplicities; and
 \item any transfer of those quadratic masses through the coherent
 physical sums requires an additional coefficient-specific theorem.
\end{enumerate}
None of the inherited inputs asserts that theorem.  In particular, the
setup supplies no affine M\"obius cancellation uniform in a moving shift,
no suppression of the full CRT face, and no parity-breaking estimate.
